% Compile with XeLaTeX or LuaLaTeX and ko.TeX. Design only; no production qualification.
\documentclass[11pt]{article}
\usepackage{kotex}
\usepackage{amsmath,amssymb,booktabs}
\usepackage[margin=1in]{geometry}
\usepackage[hidelinks]{hyperref}
\newcommand{\norm}[1]{\left\lVert #1\right\rVert}
\title{JLZ v12: 모델, 벤치마크, batch 크기를 분리한 실행 최적화 설계}
\author{실행 파이프라인 추가 명세}
\date{2026년 10월 4일}
\begin{document}
\maketitle

\paragraph{상태.}
이 문서는 구현 검토 후 작성한 실행 최적화 설계다. 새 경로의 속도나 실제 모델
수치 parity는 아직 측정하지 않았다. 기존 등록 job, source archive, method 계약을
수정하거나 실행을 재등록하지 않는다. 적용 시 별도 source/runtime 버전으로 검증한다.

\section{고정할 과학적 의미와 확인한 비용}
전체 후보 층의 subject local 목표 공동 최적화, native 문장과 tokenization,
NLL 및 KL 방향/가중치, 공유 예산, EfficiencyAdam, 요청별 moment/종료,
마지막 평가 후보, 순차 actual key 재측정, full logical-batch writer,
final-key history append와 평가 문항/분모를 유지한다.
층 수 축소, context 제거, vocabulary 근사, 반복 상한 축소, 평가 표본 축소,
새 보존 손실은 이번 실행 최적화에 포함하지 않는다.

기존 server3 ridge B1 실측은 전체 532.53초, fit 387.35초,
writer/history 76.06초, post 평가 43.53초다. W0 전체 평가 시간은 별도다.
이는 AlphaEdit writer의 시간 예측이 아니다. fit은 2,500회 physical forward와
2,400회 backward, writer capture는 600회 full forward를 기록했다.
이미 적용된 prefix 일부 재사용, microbatch별 padding 제거,
selected-position full-vocabulary head, 동일 endpoint 평가 재사용은 새 절감으로
계산하지 않는다. 누적 평가와 sequence 길이가 변하므로 B1 비중을 W20에 고정 적용하지 않는다.

\section{공통 실행 엔진과 adapter 경계}
논리 편집 batch $B$는 science 설정이며 GPU microbatch $G$와 다르다.
core는 $B=100$, 20회, 5층, width 4096/14336 또는 R/P/N 개수를 상수로 알지 않는다.
모델/벤치마크 이름으로 core에서 분기하지 않는다.
메모리에 맞춰 바꾸는 것은 $G$다. 논리 $B$를 바꾸면 joint writer의 요청 간 결합과
순차 상태가 달라지므로, batch 크기를 바꿔도 결과가 불변이라고 주장하지 않는다.

\begin{description}
\item[ModelAdapter] 층의 실행 순서, 주입/쓰기/읽기 위치, 차원, weight orientation,
native forward, mask/position 규칙, head 의미, writer 출력의 block 가법 대응,
정확한 stage stop/resume, cache dependency와 checkpoint 능력을 제공한다.
\item[BenchmarkAdapter] 순서 있는 edit occurrence, 요청별 native train/KL row와
각 reduction weight, writer key context와 별도 평균 weight, canonical row,
평가 row/후보/정답/metric/분모를 제공한다. 동일 fact의 반복 occurrence를 삭제하지 않는다.
\item[WriterAdapter] MEMIT-H 또는 AlphaEdit의 native solve, projector/history,
precision/연산 순서, 실제 weight 적용 및 history 상태를 담당한다.
\item[ExecutionPolicy] 토큰/메모리 예산, 요청 묶음, cache GPU/CPU 배치,
activation checkpoint 구간, 관측 buffering을 정한다. 목적/종료/배분 정책은 바꾸지 않는다.
\end{description}

원하는 모델 adapter가 없으면 지원을 주장하지 않는다. 학습/쓰기의 native 의미를
정의한 adapter에는 일반 full-forward 경로를 제공하고, 각 빠른 경로는 capability
검증을 통과한 경우에만 사용한다. 예를 들어 encoder의 미래 토큰은 임의로 자를 수 없고,
재사용되는 weight나 feedback 연결이 있으면 causal stage/history 재사용을 켤 수 없다.

\section{Batch 크기에 독립적인 실행 계약}
\paragraph{네 가지 크기의 분리.}
논리 편집 batch $B_t$, fit의 요청 묶음 $G_{\mathrm{fit}}$,
writer capture의 row 묶음 $G_{\mathrm{capture}}$, 평가의 row 묶음
$G_{\mathrm{eval}}$은 독립적인 설정이다. 메모리 자동 조정은 뒤의 세 크기와
토큰 예산에만 작용한다. fit의 공유 예산과 norm 계수는 요청마다 정의하고,
$B_t$나 현재 활성 요청 수에 따라 강도가 달라지도록 재정규화하지 않는다.
요청 수, context 수, target token 수는 서로 다른 축으로 관리한다.

\paragraph{요청 수와 마지막 batch.}
RunSpec은 총 편집 수 $N$과 양의 논리 batch 크기 $B$, 또는 명시적인
batch schedule을 받는다. 매 단계 실제 크기는 $B_t$이며 모든 배열, history
coverage, telemetry 분모를 이 값에서 만든다. 예를 들어 $N=2000$, $B=128$이면
128개씩 15회와 80개 한 번, 총 16회 commit이다. $B>N$이면 한 번의 $N$개
batch가 된다. 마지막 요청을 버리거나 dummy edit로 채우지 않는다.
완료 조건은 고정된 batch 번호가 아니라 정확히 $N$개의 occurrence가 반영된 것이다.
논리 batch 크기가 바뀌면 writer의 요청 간 결합과 history 갱신 시점이 바뀌므로
과학적 결과도 달라질 수 있다. 같은 논리 batch에서 GPU 묶음만 바꾼 경우에는
reference와 수치/종료 parity를 요구한다.

\paragraph{큰 batch의 writer.}
capture와 전송은 작은 묶음으로 나눌 수 있지만 solve는 전체 논리 batch에 대한
하나의 연산자다. 독립적인 소규모 write 여러 번으로 대체하지 않는다.
큰 $B$의 key/target 저장이 병목이면 다음 충분통계의 streaming을 별도 수치
backend 후보로 둔다.
\[
 C_K=KK^\top=\sum_j K_jK_j^\top,\qquad
 T_{KR}=KR^\top=\sum_j K_jR_j^\top.
\]
예를 들어 AlphaEdit의 수학적 연산자는
\[
 X=[\Pi(C_K+H)+\lambda I]^{-1}\Pi T_{KR},\qquad U=X^\top
\]
이다. 실제 구현은 역행렬을 만들지 않고 native general solve를 사용한다.
이 식은 모든 요청의 결합을 유지하지만, chunk 누적이나 $\Pi$ 곱의 재결합은
native FP32 연산 순서를 바꾼다. 따라서 별도 parity 검증 전에는 기본 경로로
켜지 않는다. layer width에 따른 행렬 메모리는 여전히 필요하다. 지원된 경로로
자원이 부족하면 필요한 자원을 보고하며, 논리 $B$를 자동으로 바꾸지 않는다.

\paragraph{평가 시점과 분모.}
평가 milestone은 누적 edit 수로 지정한다. milestone이 commit 경계와 다르면
기본적으로 해당 수를 처음 넘긴 완료 commit에서 평가하고 실제 누적 수를 표시한다.
$B=128$에서 500-edit milestone의 관측값은 512-edit 결과이며 500으로 표기하지 않는다.
정확히 500에서 평가해야 한다면 실행 전에 그 경계를 포함한 논리 batch schedule을
명시한다. GPU 편성 과정에서 임의로 batch를 나누지 않는다. 비교는 같은 실제
편집 수와 ordered prefix를 기준으로 하며, 평가 row 수는 benchmark가 제공한
실제 occurrence에서 구한다. RS/PS/NS와 ACC의 분모를 $B$, $2B$, $10B$로
고정하지 않는다.

\section{Fit: 독립 요청을 묶고 고정 계산을 옮긴다}
\subsection{요청 단위 GPU 편성}
요청마다 다른 context 수/길이를 허용한다. 한 요청의 전체 native 목적을 완성할 수
있는 row 묶음을 기본 단위로 길이별 편성한다. 매우 큰 요청은 row microbatch로
나누되 원래 context/token reduction weight로 전체 목적을 완성한 뒤 stop/backward
여부를 정한다. norm이나 공유 예산 항을 row chunk마다 중복 부과하지 않는다.
큰 논리 batch에서는 현재 활성 묶음의 변수와 요청별 optimizer state만 GPU에
둘 수 있지만, 다른 묶음의 값, moment, step counter는 그대로 보존한다.
아래 $F_r$는 native NLL과 가중 KL만을 포함하고, $\beta_r$는 v12 profile이
정한 요청별 norm 가격이다. norm gradient를 기존처럼 한 번만 더한다.
\[
 F_r=\mathrm{NLL}_r+\lambda_{\mathrm{KL}}\mathrm{KL}_r,\qquad
 L_{\mathrm{backward}}=\sum_{r\in\mathrm{continuing}} F_r,
\]
\[
 g_r=\nabla_{u_r}F_r+\beta_r\,\partial\sum_l\norm{u_{lr}}_2.
\]
request 평균, 활성 요청 수 또는 microbatch 수로 gradient를 다시 나누지 않는다.
Adam moment와 step counter는 요청별로 유지한다. 완료 요청은 변수/moment를 동결하고
다음 편성에서 제외한다. 배치화를 위해 층간 경로를 detach하지 않는다.

현재 B1처럼 100개가 모두 25회 평가한다면 $G=4$에서는 physical forward가
$25\lceil100/4\rceil=625$회, backward가 600회다. logical 평가 2,500회와
갱신 2,400회는 그대로다. 이는 호출 수 감소이며 4배 FLOPs/속도 개선을 뜻하지 않는다.
$(G,\text{checkpoint segments},\text{cache residency})$를 함께 메모리/시간으로 선택한다.
평가 성능으로 편성을 선택하지 않는다. OOM 재편성은 미완료 candidate의 gradient와
일시 상태를 버리고 같은 candidate를 재실행하며, 요청/토큰을 누락하지 않는다.
요청별 update는 필요한 계산 완료 후 게시하고, 재시도에서 moment/step을 두 번
갱신하지 않는다. 재편성에 따른 row 순서는 native reduction 순서와 구분해 기록한다.

\subsection{첫 주입 직전 출력 cache}
fit에서는 weight가 고정이고 첫 delta가 block 출력에 더해진다. 따라서 해당
block의 native 출력 전체를 entry에서 한 번 저장한 뒤 $h^0+\delta$부터 실행할 수 있다.
현재 구현은 first key와 residual을 CPU에 저장하고 candidate마다 옮겨
첫 down projection을 다시 계산한다. Llama profile에서 두 feature의 폭
$14336+4096$ 대신 $4096$을 저장하면 그 feature 부분은 4.5배 작아진다.
mask/position 등 나머지 cache와 전체 메모리/속도에 대한 배율은 아니다.
cache는 첫 주입보다 뒤의 delta 의존 activation을 포함하지 않는다.
첫 주입 경계는 현재 비제로 층이 아니라 전체 eligible 변수 중 가장 앞선 위치다.
0인 층의 gradient와 재진입을 보장하기 위해 경계를 뒤로 옮기지 않는다.
첫 block weight가 다음 edit batch에서 바뀌므로 이 output cache는 batch entry에 결속한다.

\subsection{GPU 상태와 telemetry}
KL teacher, row owner/lookup/scoring 위치와 가능한 entry cache를 RAM 예산 안에서
GPU에 상주시킨다. AlphaEdit projector는 fit에 쓰이지 않으므로 writer 단계로
늦춰 로드하고 한 층씩 관리할 수 있다. 현재 다섯 projector는 FP32로 약 3.83GiB다.
이는 CPU/전송/메모리 사용을 재배치하는 선택이지 수학 변경이 아니다.

canonical hidden은 같은 forward에서 capture하되, 계속 학습하는 후보에서는
GPU에서 필요한 scalar 진단만 계산한다. terminal로 결정된 요청만 실제 $z_l$를
복사한다. GPU scalar를 요청/층마다 Python float로 읽지 않고 묶어서 전송한다.
Optimizer는 같은 요청별 연산을 tensor화하며, 모든 후보 층의 gamma와 재진입을 유지한다.
이벤트 내용과 candidate identity는 유지하고 writer handle을 열어 둔 채 bounded buffer로
기록한다. round/terminal/transaction 경계에서 flush하며 commit은 필수 이벤트의
영속화 이후에만 게시한다. crash 시 미flush 구간은 완성 기록으로 주장하지 않는다.

\section{Writer: 실제 key 재측정은 유지하고 불필요한 경로를 없앤다}
\subsection{첫 적용: 층까지만 capture}
현재 capture는 한 층의 key/hidden을 요청해도 모델 끝까지 돈다. native forward의
요청된 마지막 지점에서 종료하도록 한다. canonical block 출력이 필요하므로
key 입력만 얻고 그 block 전부를 건너뛰지는 않는다.
현재 32-block 모델과 L4--L8의 다섯 prewrite 및 한 번의 final capture에서는
동일 row group 기준 block 실행 수가 $6\cdot32=192$에서
$5+6+7+8+9+9=44$로 줄어든다. layer 비용이 같다는 근사 아래의 capture 작업량 비교이며,
solve/history/관측/fit을 포함한 전체 속도 배율이 아니다.
Writer용 row는 benchmark adapter가 정의한 native key 문장만 사용하고
학습용 KL row를 불필요하게 실행하지 않는다.

\subsection{후속 적용: 층 순서대로 멈추고 재개}
whole logical batch의 native key와 canonical hidden을 층 $l$까지 streaming capture한다.
전체 $K_l,R_l=z_l^v-h_l^{cur}$가 모인 뒤 solve를 한 번 수행하고, 실제 저장된
새 weight로 그 층 출력을 만든 뒤 다음 층으로 진행한다. 메모리가 부족하면
stage state를 bounded CPU buffer로 옮긴다. 다음 층의 key는 갱신된 하층 출력에서
새로 계산되므로 actual-key 계약이 유지된다. microbatch별 독립 writer로 나누지 않는다.

실제 출력을 $h_{before}+U k$로 임의 대체하지 않는다. 원래 residual/linear/cast/add
순서 또는 검증된 native block 재실행으로 계산하여 FP32 차이를 검증한다.

\subsection{Final-key 재사용의 조건}
layer-ordered feedforward 모델에서 층 $l$의 key가 자기 변경 weight보다 앞에 있고,
이미 반영한 하층 외에 그 key로 되돌아오는 의존성이 없다면
\[
 K_l^{\mathrm{before\ own\ write}}=K_l^{\mathrm{final}}.
\]
후속 상층 write는 하층 key를 바꾸지 않는다. 이 조건을 adapter가 입증하고 native
context/token 규칙과 parity가 맞으면 write 직전 key를 저장해 둘 수 있다.
history는 여전히 모든 write 성공 뒤 원래 native dtype와 순서로 한 번 append한다.
최종 hidden 관측도 실제 갱신 block에서 capture한다. 조건이 없으면 기존 final
forward를 유지한다. 원래 문서의 ``필수 final forward''는 이 capability가 있는
경우 ``필수 final-state key 값''으로만 대체하는 명시적 실행 계약 변경이다.

\section{평가와 writer 선형대수}
\paragraph{평가.}
동일 모델 상태, token IDs, mask/position, readout을 갖는 forward 입력을 공유한다.
원 평가 occurrence와 target별 채점/분모는 모두 유지한다. 현재 profile처럼 target이
한 token이면 new/true의 입력이 같은 경우가 많아 같은 hidden/logits로 두 target을
채점할 수 있다. 다중 token 후보는 전체 입력이 다르면 자동 공유하지 않는다.
공통 prefix KV 공유는 별도 causal capability와 parity가 있을 때만 적용한다.
W0 cache는 method 이름 대신 실제 model/input/tokenizer/evaluator/runtime 정체성으로
재사용 여부를 판단한다. 다른 endpoint의 logits/activation을 그대로 쓰지 않는다.

\paragraph{AlphaEdit solve.}
최초 최적화에서는 native FP32의 일반 선형계와 연산 순서를 그대로 둔다.
projected basis, Woodbury 또는 history rank-update factorization은 수학적으로
검토할 수 있지만, FP32 projector의 근사 idempotence와 history 누적 오차, 연산
순서가 달라진다. 조건수와 applied update/후속 성능을 검증한 별도 numeric backend로
분리한다. 단순 batch/width 비율을 전체 가속 배율로 주장하지 않는다.

\section{적용 순서와 확인 기준}
1차는 portable interface와 원 reference executor, 전체 요청 GPU 묶음, terminal-only
복사/진단 buffering, phase별 cache residency다. 2차는 first-injection cache,
writer capture 조기 종료, 평가 입력 공유다. 3차는 adapter가 지원하는 staged writer와
final-key 재사용이다. writer 대수 변경/새 precision/attention kernel은 별도 후순위다.

논리 batch와 microbatch를 독립적으로 변화시키고, 마지막 partial batch,
다른 context 수/target 길이, 비연속 후보 층, 가능한 heterogeneous width를 검사한다.
작은 qualification에서는 $B\in\{1,3,7\}$과 서로 다른 fit/capture/eval 묶음을
사용하고, $N=10,B=3$의 마지막 한 요청 및 $B>N$을 포함한다.
2000-edit schedule은 $B=100$의 20회와 $B=128$의 16회/마지막 80개를
CPU에서 검산한다. 이 schedule 검산은 모델 parity나 속도 검증을 대체하지 않는다.
작은 모델/요청에서 reference와 후보별 NLL/KL/gradient, 요청별 stop/moment,
terminal $z$, applied $U$, 현재/최종 key, history와 raw metric을 비교한다.
decision boundary 근처 차이는 원 경로 재평가로 판정하고 원인을 기록한다.
기존 forward/gradient tolerance를 임의로 완화해 PASS하지 않는다.
새 모델/벤치마크 조합은 각 adapter qualification을 별도로 받는다.

짧은 동일 workload 측정은 fit, transfers, optimizer/telemetry, key capture,
solve, history, evaluation, state verification을 분리한다. useful/padded tokens,
logical 후보 수, physical forward/backward 및 checkpoint 재계산,
GPU peak/CPU RSS를 함께 보고한다. 전체 시간은 각 단계의 실측을 합산하며
미측정 가속 배율이나 AlphaEdit 전체 ETA를 확정하지 않는다.
\end{document}
